\begin{abstract}

Long-distance quantum networks require repeaters that distribute high-fidelity Bell pairs despite photon loss and operational noise. In all-photonic quantum repeaters built from \Gls{hrgs-long}, entanglement purification can be applied at different stages, from resource generation to individual hops and end-to-end entanglement. The number of copies and circuit used may also vary between stages, making the choice of where and how to purify a combinatorial scheduling problem. The source papers demonstrate the potential of such scheduling through a single hand-picked example, but do not systematically explore the space of possible schedules.

This thesis represents each fixed non-adaptive purification schedule as a \Gls{dag} and evaluates it with a single bottom-up pass computing fidelity, success probability, latency, rate, and generation cost. The evaluator reproduces the source paper's three canonical fidelity benchmarks to four significant figures. Three search tiers then construct and compare candidate schedules under a shared generation-cost budget $E_{\max}$: structured baseline enumeration, an exact span-partition dynamic program, and beam search with an additional same-span pumping operation. All candidates are structurally validated and evaluated using the same physical model throughout the protocol.

At the source papers' reference configuration, a schedule found by search at the same generation cost as the hand-picked baseline improves the output rate by $2.5\%$. With a relaxed cost constraint, search finds a schedule requiring only $60\%$ of the baseline generation cost while achieving a $52.8\%$ higher output rate. More importantly, when inner-qubit error is introduced, the fixed baseline can fall below the required fidelity, whereas a schedule can be found by search at the same generation cost while remaining feasible. The same feasibility advantage is observed on average across independently randomized per-hop heterogeneous networks. Further experiments examine its dependence on noise level, timing, and chain length.

These results establish existence within the implemented schedule families and search limits and not global optimality: beam search and its pumping operation retain only a bounded frontier at each stage and are therefore only heuristic. Consequently, failing to find a feasible schedule does not rule out its existence, as it may lie outside the explored search frontier.

\glsreset{hrgs-long}
\glsreset{dag}

\end{abstract}